\section{Metoda analizy wykonania martwego ciągu}

\subsection{Przebieg algorytmu analizy}

Algorytm analizy wykorzystuje parametry wyznaczone dla kolejnych klatek nagrania i na ich podstawie określa przebieg wykonywanego ruchu. Analiza prowadzona jest w czasie, dzięki czemu możliwe jest wydzielenie pojedynczych powtórzeń.

Dla każdego wykrytego powtórzenia wyznaczany jest zestaw cech opisujących jego przebieg, między innymi zakres ruchu, czas trwania oraz wartości wybranych parametrów w określonych momentach ćwiczenia. Dane te stanowią następnie podstawę regułowej oceny powtórzenia.

Po przeanalizowaniu poszczególnych powtórzeń ich parametry są dodatkowo porównywane w celu oceny spójności całej serii. Szczegółowe zasady wykrywania i oceny powtórzeń przedstawiono w kolejnych podrozdziałach.

\subsection{Wygładzanie przebiegu ruchu}

Przed wykrywaniem powtórzeń wygładzany jest przebieg pionowego położenia sztangi. W tym celu zastosowano średnią kroczącą z oknem pięciu próbek. Dla próbek, dla których dostępne jest pełne okno, można ją zapisać jako:

\begin{equation}
\hat{y}_i =
\frac{1}{5}
\sum_{k=-2}^{2} y_{i+k},
\end{equation}

gdzie \(y_i\) oznacza pionowe położenie sztangi dla kolejnej analizowanej klatki, a \(\hat{y}_i\) wartość po wygładzeniu. Na początku i końcu przebiegu, gdzie pełne okno nie jest dostępne, średnia wyznaczana jest na podstawie dostępnych próbek.

Wygładzony przebieg jest następnie wykorzystywany podczas wykrywania powtórzeń, co ogranicza wpływ chwilowych zmian położenia wyznaczanego na podstawie kolejnych klatek.

\subsection{Identyfikacja faz ruchu i wydzielanie powtórzeń}

Wykrywanie powtórzeń opiera się na zmianach wygładzonego pionowego położenia sztangi. W przyjętym układzie współrzędnych większa wartość położenia oznacza niższe położenie sztangi. Na tej podstawie lokalne maksima przebiegu odpowiadają dolnym pozycjom ruchu, natomiast lokalne minima jego górnym pozycjom.

Za pojedyncze powtórzenie uznawana jest sekwencja dolna pozycja--górna pozycja--dolna pozycja, spełniająca przyjęte warunki dotyczące zakresu i czasu trwania ruchu. Początek fazy podnoszenia wyznaczany jest po odpowiedniej zmianie położenia sztangi względem wykrytej pozycji dolnej. Fragment bezpośrednio poprzedzający podnoszenie może zostać zakwalifikowany jako faza przygotowania.

Dla wykrytego ruchu wyróżniane są fazy przygotowania, podnoszenia i opuszczania, natomiast pozostałe fragmenty nagrania oznaczane są jako stan spoczynku. Wyznaczone granice powtórzeń są następnie wykorzystywane do obliczenia ich charakterystyk i dalszej oceny.

\subsection{Wyznaczanie charakterystyki pojedynczego powtórzenia}

Po wyznaczeniu granic powtórzenia obliczane są parametry opisujące jego przebieg. Określany jest czas trwania ruchu oraz zakres pionowego przemieszczenia sztangi. Dodatkowo zapisywane są wartości kątów stawu kolanowego i biodrowego w wybranych momentach ruchu oraz zmiana nachylenia pleców w trakcie powtórzenia.

Dla całego powtórzenia wyznaczana jest również zmiana kąta nachylenia pleców. W pozycji początkowej porównywane jest położenie barku i biodra. Tak przygotowany zestaw cech stanowi dane wejściowe dla reguł oceny techniki opisanych w kolejnym podrozdziale.

\subsection{Regułowa ocena pojedynczego powtórzenia}

Ocena pojedynczego powtórzenia wykonywana jest na podstawie wcześniej wyznaczonej charakterystyki ruchu. Sprawdzany jest między innymi osiągnięty zakres pionowego przemieszczenia sztangi, porównywany z przebiegiem pozostałych powtórzeń w analizowanej serii. Pozwala to wskazać wykonania o ograniczonym zakresie ruchu.

Oceniana jest również pozycja końcowa na podstawie wartości kąta kolana oraz przybliżonego kąta pomiędzy tułowiem a udem. Parametry te pozwalają określić stopień wyprostu osiągnięty w górnej pozycji ruchu. Dodatkowo analizowane jest ustawienie tułowia w pozycji początkowej oraz zmiana jego nachylenia w trakcie powtórzenia.

Określana jest także charakterystyka pozycji początkowej. Na podstawie wartości kąta kolana może ona zostać opisana jako niższa lub wyższa względem typowej pozycji obserwowanej w analizowanej serii. Informacja ta ma charakter opisowy i nie jest sama w sobie traktowana jako błąd wykonania.

\subsection{Ocena spójności serii}

Po ocenie pojedynczych powtórzeń ich parametry są porównywane w celu określenia spójności wykonania całej serii. Analizowany jest zakres pionowego ruchu, czas trwania powtórzeń oraz wartości kątów stawu biodrowego i kolanowego w górnej pozycji ruchu.

Dla każdego z tych parametrów wyznaczana jest różnica pomiędzy największą i najmniejszą wartością uzyskaną w serii. Przekroczenie przyjętego dopuszczalnego zakresu zmienności powoduje wskazanie braku spójności danego elementu wykonania. Ocena ta jest wykonywana oddzielnie od oceny techniki poszczególnych powtórzeń i opisuje powtarzalność ruchu w obrębie całej serii.
